%
%	WISS 2026 本文ドラフト
%	結果（4.5-4.8）・5.1（結果の解釈）・6（まとめ）・概要は
%	実験実施後に追記する。書式は wiss_template.tex に準拠。
%
\documentclass[twoside]{wiss}

\usepackage[dvipdfmx]{graphicx}
\usepackage{color}
\usepackage{flushend}
\usepackage{url}

\journalhead{運動準備電位の先読みによる想起オンセット同期（仮）} %%%%%% ←←　タイトル確定後に修正すること

\begin{document}

\title{運動準備電位の先読みによる想起オンセット同期\\
――運動イメージと仮想手の動き出しのラグをなくす VR システム\textcolor{red}{（仮タイトル）}}
\etitle{}

\author{（著者未定）}

\begin{abstract}
% 概要（400字程度）は実験結果確定後に執筆する。
% WISS の慣例に従い「概要．」と本文を改行せず一続きで書くこと。
\end{abstract}

\maketitle

\section{はじめに}

ブレインコンピュータインターフェース（BCI）は，脳活動を計測し，その情報をもとに外部機器を操作する技術である．筋肉や末梢神経を経由せずに，脳の活動そのものからコマンドを取り出せることから，随意運動を介さずに意図を機器へ伝えられる．非侵襲な脳波（EEG）計測の普及とともに，医療・リハビリテーションの領域を中心に発展してきた BCI は，近年ではこうした医療用途にとどまらず，VR/HCI 分野など，健常なユーザを対象とした操作モダリティとしても応用が広がっている．身体を動かさずに機器を操作できるという特性は，ゲームや VR のように身体表現そのものが体験の中心となる領域において，新しい操作感覚をもたらす可能性がある．

特に運動イメージ（MI）に基づく BCI は，身体を動かす想像に伴う感覚運動リズムの変化を特徴量とするパラダイムであり，脳波の変化が身体運動と対応しやすいため，他の BCI パラダイムに比べて操作の意味づけを理解しやすい．実際に手を動かす代わりに，動かす様子を思い浮かべるだけでよいという操作方法は習得の負担が小さく，VR 空間での身体表現とも相性がよいことから，広く用いられている．本研究が対象とする，仮想手を操作して物を握るという課題も，この運動イメージパラダイムの延長線上にある．

しかし MI-BCI では，脳波の変化から意図を推定するために一定の解析時間窓を必要とするため，イメージ想起の開始から仮想的な動作として反映されるまでに，数百ミリ秒規模の遅延が避けられない．この遅延は，単なる操作の不便さにとどまらず，身体所有感や操作感そのものを損なう．とりわけ，仮想手を自分の身体の延長として提示する VR システムでは，想起と動作のずれがそのまま身体表現の破綻として体験されやすい．遅延を挿入したロボット手の錯覚実験では，実際の手の動きと仮想手の動きの時間的なずれが 190ms を超えると，身体所有感が消失することが報告されている\cite{ismail2016}．

この遅延をなくす試みの多くは，特徴量抽出に用いる時間窓を短くしたり，計算量の少ないフィルタを用いたりすることで，信号処理そのものを高速化する方向で進められてきた．しかし，遅延の主要因は解析時間窓の長さやフィルタの群遅延にあり，計算処理を高速化するだけでは越えられない．加えて，こうした対処の多くはイメージ想起の開始，すなわち動き出しの瞬間を早く検出することに主眼を置いており，動き出した後の操作の維持まで一貫して遅延をなくす構成は十分に確立していない．動き出しの検出だけを速くしても，その後の把持や保持の操作に遅延が残れば，体験全体としての遅延は解消されない．

そこで本研究では，想起開始の前後で役割を分けた2つの分類器を組み合わせることで，運動想起の開始と仮想手の動き出しを同期させたうえで，その後の操作の維持までを一貫して扱う VR システムを提案する．一方の分類器は運動開始に先行する脳活動から動き出しの瞬間を検出し，他方の分類器はその後の操作の継続と終了を判別する．動き出しの同期と操作の維持という異なる要求を単一の分類器に担わせるのではなく，役割ごとに専用の分類器を割り当てることで，双方の要求を無理なく両立させることをねらう．用いる具体的な信号については3章で述べる．

\section{関連研究}

\subsection{運動関連の皮質活動による運動前検出}

事象関連脱同期（ERD）は，運動の準備・実行・イメージに伴って，感覚運動野における $\mu$/$\beta$ 帯律動のパワーが低下する現象であり，運動イメージに基づく BCI で標準的に用いられる特徴量である．右手の運動イメージであれば左半球で，左手であれば右半球でパワー低下が強く現れるという対側優位性を示すことから，左右どちらの身体部位を想起しているかの判別にも利用できる．この対側性は，ERD が単なる覚醒度の変化ではなく，運動を司る脳部位に固有の活動を反映していることの根拠でもある．

ただし ERD は，一定の時間窓内での帯域パワー，すなわち信号の分散を推定することによって得られる特徴量であるため，その定義上，運動やイメージの開始後の時間窓を要する．したがって，運動開始に先行して検出することは原理的にできない．一方，運動準備電位（MRCP）は 0.05〜5Hz の低周波帯に現れる緩電位であり，ERD が用いる $\mu$/$\beta$ 帯とは周波数帯域が重ならない．このため，同一の脳波信号から，両者を互いに干渉させることなく独立に抽出できる．

MRCP は，随意運動の約1.5〜2秒前から緩やかに始まり，運動開始が近づくにつれて傾きが急峻になりながら陰性方向に振れ，運動開始の直前から直後にかけて陰性のピークに達する電位である．中心部電極（C3・Cz・C4）を中心に，振幅数マイクロボルトの緩やかな変化として観測され，運動が始まる前に運動の準備を検出できる唯一の生理的特徴量である．

単一試行での MRCP 検出手法は，空間フィルタから部分空間学習，信号源分離へと世代を重ね，検出の潜時を縮めてきた．第1世代の空間フィルタを用いた手法は，単一試行データに対するオフライン評価において，検出率63.5\%，運動開始の約203ms前という，運動開始に先行する検出を報告した\cite{niazi2011}．続く第2世代は，部分空間学習（局所性保存射影とその後段の線形判別分析）を用いてオンライン環境での検出を初めて評価した\cite{xu2014}．さらに第3世代として，制約付き独立成分分析による信号源分離を用いた手法も提案されている\cite{karimi2017}．一方，近年ではアルファ律動と MRCP という周波数帯の異なる2種類の特徴量をそれぞれ独立に検出したうえで決定レベルで融合し，検出精度と待ち時間の両立を図る手法も報告されている\cite{liu2024}．検出手法が世代を追うごとに複雑化しているのは，単一試行かつオンラインという制約のもとで検出精度を高めるための代償であるといえる．

ただし，こうした検出手法によって運動開始より前に安定して検出することは，オンライン環境では依然として難しい．第2世代の手法によるオンライン評価でも，検出潜時は平均315msにとどまっており，これは運動開始に先行する値ではなく運動開始後の値である\cite{xu2014}．すなわち，検出手法が世代を重ねてもなお，検出は運動開始の前ではなく後になされているのが実情である．この差は，オフライン解析では試行全体の波形を前後から参照できるのに対し，オンライン処理ではその時点までの情報しか使えないという制約の違いに起因すると考えられる．

\subsection{BCI における遅延の影響}

実身体を伴う操作では，遅延はまず運動パフォーマンスと同時性知覚という2つの側面を，75ms という早い段階から損なう．すなわち，操作の正確さや速さそのものが低下し始めると同時に，自分の動作と画面上の結果が同時に起きたと感じられるかどうかの判断も曖昧になり始める\cite{waltemate2016}．こうした低下は，遅延が明確に知覚できるほど大きくなるよりも前の，比較的軽微な遅延の段階からすでに生じている点で見過ごされやすい．

遅延がさらに大きくなると，行為主体感や身体所有感といった身体化の感覚も 125ms から低下し始め，300ms を超えると顕著に劣化する\cite{waltemate2016}．遅延を挿入したロボット手の錯覚実験でも同様の傾向がみられ，身体所有感は 190ms を超える遅延で消失する一方，行為主体感は 290〜490ms の遅延条件でも部分的に残存することが示されている\cite{ismail2016}．身体所有感と行為主体感とでは，遅延に対する頑健性そのものが異なるといえる．

一方で，BCI を介した操作については，これとは相反する報告もある．運動想起によって駆動される脳機械インターフェース操作を対象とした研究では，脳活動と視覚フィードバックのあいだの遅延が1秒を下回る範囲では主体感の明確な低下が見られなかった．これは，実身体の動きを伴わない BCI 操作では，主体感の判断が身体感覚ではなく視覚フィードバックにほぼ支配されるためだと説明されている\cite{evans2015}．裏を返せば，視覚フィードバックさえ矛盾なく提示できれば，脳波解析側の遅延はある程度許容されうることを意味する．

この遅延の許容度の違いは，許容が固定された値ではなく，操作の文脈によって変わりうることを示唆する．一人称視点の VR アバターによる身体化の程度と BCI 操作の遂行度は，二値運動イメージタスクにおける行為主体感の分散のうち41\%を説明することが報告されている\cite{ziadeh2021}．これは，遅延量そのものよりも，どれだけ自分の身体として感じられているかという身体化の程度が，主体感を左右する一例といえる．

ただし，遅延したニューロフィードバックは，単に主体感を損なうだけでなく，学習そのものを阻害することが報告されている\cite{belinskaia2020}．フィードバックが遅れて届くと，脳活動の変化とフィードバックの対応づけ自体が難しくなり，訓練が成立しなくなるためである．したがって，BCI の較正・訓練フェーズでは，操作時以上に遅延を減らすことの優先度が高くなる．本システムが起動判定の閾値調整を提示側だけで完結させているのも，この訓練フェーズでの遅延を最小限に抑えるための設計である．

\section{システム}

\subsection{システム構成}

本システムは，脳波の計測を EMOTIV EPOC Flex 2 で，仮想手の提示と制御を Meta Quest 3 上のアプリケーションで行う．脳波は32チャネルのゲル電極を用い，サンプリング周波数256Hzで計測する．システム全体は，脳波を解析して分類結果と確信度を出力する側と，その結果を受けて仮想手を制御する提示側とに分かれており，いつ手を動かし始めるかという判定は提示側に委ねる設計とした．脳波解析側は，運動準備電位や事象関連脱同期といった特徴量から分類結果と確信度を計算するところまでを担い，その値をどう解釈して仮想手を動かすかという判断には関与しない．こうして判定条件を脳波解析側の実装から切り離すことで，被験者ごとに必要な閾値や条件の調整を，提示側の設定だけで完結できるようにしている．この分離は，脳波解析側のアルゴリズムを変更しても提示側の較正手順に影響が及ばないという意味で，判定条件を明示したまま追試できる実験基盤としての性格も持つ．

\subsection{信号処理・特徴量}

本システムは，運動開始前から生じる運動準備電位を主な特徴量として用いる．2.1節で述べたとおり，生理的な遅延そのものは信号処理の高速化によっては削れず，運動開始に先行する信号を先読みすることによってしか越えられない．運動準備電位は，先読みによって体感的なラグを負の値，すなわち動作の想起よりも先に仮想手が動き出す状態にまで縮められる，数少ない生理的経路である．事象関連脱同期のように運動開始後の情報に依存する特徴量では，原理的にこの負のラグを実現できない．

脳波はリアルタイム推定に適した計測系で取得する．運動野を覆う中心部電極（FCz・Cz・C3・C4）を中心に計測し，眼球運動由来のアーチファクトは前頭極電極（Fp1・Fp2）を用いて除去する．前処理には 0.1〜3Hz の因果的な帯域フィルタを適用する．波形の位相ひずみを避けるためによく用いられる零位相フィルタは，処理対象より未来の値を必要とするためリアルタイム処理には使えず，本システムでは過去の値のみから逐次計算できるフィルタを用いる．因果的フィルタは零位相フィルタに比べて位相のずれを伴うが，本システムではリアルタイム性を優先し，この位相のずれは後段の判定側で織り込んだうえで扱う．

この生波形から，運動準備電位はキュー直前の時間窓におけるベースライン補正後の振幅を，事象関連脱同期は $\mu$/$\beta$ 帯のバンドパワーを，それぞれ独立に特徴量として抽出する．両者は周波数帯域が重ならないため，同一の脳波から互いに干渉することなく取り出せる．先行研究では2つの特徴量を決定レベルで融合して検出精度そのものを高めるのに対し\cite{liu2024}，本システムでは両者を起動と維持という異なる役割に割り当てる点が異なる．目的に応じて特徴量の使い道を分けることで，それぞれの特徴量が持つ強みを活かしたまま，性質の異なる2つの判定を同時に成立させることをねらう．

\subsection{起動・維持の分類}

本システムは，手を握り始める起動の判定と，握りを保ち続ける維持の判定を，別々の分類器に分ける役割分離のアンサンブルを核とする．運動準備電位は，運動開始の前後に限って現れる一過性の電位であり，握りを持続的に駆動する信号としては使えない．そのため，起動時の一過性の信号と，維持中に持続する信号とで，役割そのものを分ける必要がある．単一の分類器にすべてを担わせようとすると，性質の異なる2種類の入力を同じ枠組みで扱うことになり，どちらの検出も中途半端になりかねない．

起動は，運動準備電位から安静と握るイメージの開始を判別する二値の分類器が担う．この分類器が握るイメージの開始を検出した時点で，仮想手の動き出しが引き起こされる．この分類器の役割は，安静状態から握るイメージへの遷移という一度きりのイベントを検出することに限定される．

起動後の維持は，事象関連脱同期から握るイメージの継続と安静への復帰を判別する別の分類器が担う．この分類器の出力に応じて，握りの保持と解放が切り替わる．起動用の分類器とは異なり，この分類器は試行中に出力を更新し続けることで，握り続けている間の状態を継続的に追跡する．

いずれの分類器も，較正データが少なく実時間制約に適合する軽量な線形分類器として，正則化線形判別分析（shrinkage LDA）を用いる．1セッションあたりの較正試行数は限られており，深層学習のような大量のデータを要する手法よりも，少ないデータから安定して学習できる古典的な線形手法のほうが適している．shrinkage LDA は，共分散行列の推定を正則化することで，少数の較正試行からでも過学習を抑えた判別境界を求められる点が，この制約に合致する．

\subsection{調整目標}

運動イメージは，実際の身体運動とは異なり，物理的な開始時刻を持たない．そのため，動作の一致を筋電などの物理指標だけで測ることはできない．そこで本システムは，キューを基準とした起動時刻の誤差に加えて，被験者自身が回答する主観的なタイミング一致度を対にして記録する．客観的な誤差だけでは，被験者が実際にどう感じたかという体験の質を捉えられないためである．

そのうえで本システムは，この主観的な一致度を，起動判定を行う分類器を調整するための目標として用いる．主観評価を，結果を記述するための指標としてのみ扱う既存研究に対し，主観をデコーダ調整という設計プロセスの内部に組み込む点が，本システムの新規性の核である．客観的な起動誤差と主観的な一致感を対にして扱うことで，デコーダの調整が，数値上の精度だけでなく，被験者にとっての体験の質にも向かうようにする．

\subsection{VR提示}

提示側は，分類器の確信度が閾値を超える状態が一定回数連続したときに仮想手の動き出しを起こし，その直後に不応期を設けることで，単発の跳ねに左右されず再現可能な起動条件を保つ．閾値と連続回数は，各被験者の較正データを用いて個別に決定し，本番のタスク中は固定する．こうして判定条件を外から見える形にすることで，条件を明示したまま追試できる構成とした．閾値や連続回数を本番中に動的に変更しないことで，セッション間・被験者間の比較においても，起動条件そのものが交絡要因とならないようにしている．

仮想手は一人称視点で提示し，起動判定に応じて開閉させる．運動開始の約2秒前から始まる運動準備電位の陰性化を，運動開始に先行する時点で検出できた場合には，想起の開始と仮想手の動き出しが同期し，体感的なラグが生じない．検出が運動開始よりも後になった場合には，一般的な BCI と同様に，一定の遅延を伴う同期にとどまる．

開閉動作は一定速度に制限し，仮想手が把持に到達する時刻を計測できるようにする．試行ごとに，起動時刻と被験者による主観評価をあわせて記録する．速度を固定するのは，把持到達時刻のばらつきが，起動判定の精度に由来する変動と，提示側の描画に由来する変動とで混ざらないようにするためである．

\section{ユーザスタディ}

\subsection{参加者}

本研究の評価は，本システムが複数日にわたる運用に耐えることを確かめる，単一被験者・複数日のユーザスタディとして行う．評価の目的は多日運用の実証であり，多数の被験者を対象とした統計的な母集団推定は目的としない．BCI は，電極位置のわずかなずれや被験者の覚醒度など，日ごとの非定常性の影響を受けやすい．そのため，同一の被験者を対象に複数日にわたって計測を重ねることで，こうした非定常性のもとでも本システムが機能し続けるかを確認する．

\subsection{手順}

各日のセッションは，装着と較正，トレーニング計測，モデル構築，タスク，事後アンケートという同一の手順を，複数日にわたって反復する．装着と較正では，電極位置と信号品質をその日ごとに確認したうえで，当日の状態に合わせて分類器を再構築する．同一の手順を繰り返すことで，日をまたいだ性能や主観評価の推移を追跡できるようにする．

\subsection{タスク}

各試行は，安静，予告，キュー，保持，復帰という構造で進行する．安静区間には，周期的な起動との区別のためわずかな時間の揺らぎを与える．予告からキューまでの間隔を不規則にすると運動前の皮質活動そのものが減弱することが報告されているため\cite{cui2009}，本システムでは予告からキューまでの間隔を固定する．間隔を固定することは，被験者がキューの到来をある程度予期できることを意味するが，これは運動準備電位そのものを安定して誘発するために必要な設計である．

トレーニングとタスクは同一の試行構造を共有し，起動の有無と主観回答の有無だけが異なるようにする．こうすることで，トレーニングで学習したモデルが想定する入力の分布を，タスク中に大きくずらさないようにする．

各試行の直後には，被験者は早すぎから遅すぎまでの7段階で，主観的なタイミング一致度を回答する．

\subsection{評価指標と解析}

評価指標は，キューを基準とした起動時刻の誤差を主要指標とし，あわせて検出率，誤起動率，把持到達時刻の誤差，試行成功率，主観的一致度，行為主体感，身体所有感，作業負荷，シミュレータ酔いを日ごとに記録する．起動しなかった試行は検出率の分母には含める一方，起動時刻誤差の分布からは除外することで，検出できなかった試行の存在と，検出できた試行の精度とを混同しないようにする．作業負荷とシミュレータ酔いは，複数日にわたる運用の負担そのものを評価するために，タスクの精度指標とは別に日ごとに記録する．

解析では，起動時刻誤差の分布，起動時刻誤差と主観的一致度の被験者内相関，日およびブロック間の推移，脳波に依存しない対照条件との比較を検討する．対照条件との比較は，観測された起動が偶然のタイミングの一致ではなく，脳波の変化そのものに由来することを確かめるために設ける．

\section{議論と限界}

% 5.1 結果の解釈は実験結果確定後に追加する。
% 追加すると本節以降の \subsection は自動的に 5.2, 5.3 に繰り下がる。

\subsection{限界}

運動イメージでの検出率は，実際に身体を動かす運動実行に比べて低くとどまることが知られている．本システムも，この特性を前提としたうえで，一定の割合の試行では起動しないことを許容する設計とした．検出できなかった試行を無理に起動させようとせず，検出できなかったこと自体を記録する設計は，見かけ上の成功率を高く見せることよりも，システムの実態を正確に記述することを優先した結果である．

また，キューを基準とした起動時刻誤差には，被験者自身の想起開始そのもののばらつきに由来するノイズが含まれる．そのため，起動時刻誤差として観測される値は，システムの検出精度だけでなく，被験者の想起のばらつきも反映しており，真値そのものに不確実性がある．

加えて，本実験の構成では，運動準備に由来する成分と，予告への時間的な期待に由来する随伴陰性変動（CNV）成分とを分離することができない．いずれも中心部電極上の緩やかな陰性電位として現れるため，運動準備電位単体を切り出すことは容易ではない．この混入は，分類器が運動準備そのものではなく，予告に対する時間的な期待を検出している可能性を排除できないことを意味する．

運動準備電位の解釈自体にも，運動の意識的な開始点を反映するのか，それとも自発的な神経活動のゆらぎが閾値を超えた結果を反映するに過ぎないのかという論争がある\cite{schurger2012}．そのため本研究は，意図の先読みそのものを主張するのではなく，運動準備に伴う皮質活動を工学的に検出することにとどめる．

さらに，本実験の対照条件は，脳波と無関係なタイミングで仮想手を動かす条件を含んでいない．そのため，仮想手の起動が脳波に随伴していること自体がもたらす効果までは，本実験の範囲では示せない．この点を確かめるには，脳波と無関係な周期的タイミングで仮想手を起動させる対照条件を別途設ける必要がある．

\subsection{今後の展望}

取り除けない残余の遅延は，制御対象と提示対象の変位比を操作することで重さの知覚を誘発する疑似触覚の技法を用いることで\cite{samad2019}，システムの遅さそのものではなく，物体の重さとして意味づけ直す拡張が考えられる．運動準備電位による同期の上にこの技法を重ねる構成は，まだ前例がない．遅延を隠すのではなく，遅延の存在を物体の性質として体験に組み込み直すという発想は，本システムが目指す遅延の解消とは異なる角度からの補完になりうる．

また，本研究は単一被験者によるケーススタディであるため，より多人数での検証へ拡張し，統計的な一般化可能性を確認することが今後の課題である．

%%
%%	参考文献
%%
\bibliographystyle{jwiss}
\bibliography{references}

\end{document}
